\section{Décomposition spectrale bicouche de l'opérateur angulaire}

Soit \(R=R^*=R^{-1}\) l'involution active et

\[
P_\pm=\frac{I\pm R}{2},
\qquad
T=q_+P_+ +q_-P_-,
\qquad 0<q_+<q_-<1.
\]

Avec
\[
x=A_{\rm rad}=\frac{\pi A}{180},\qquad
y=C^*_{\rm rad}=\frac{\pi C^*}{180},
\]
les deux canaux sont construits avant d'évaluer le vide :
\begin{equation}
q_+=e^{-x-y},\qquad q_-=e^{-x+y}.
\label{eq:vacuum-angular-channels}
\end{equation}
L'échange de feuillets inverse \(C^*\), c'est-à-dire \(y\mapsto-y\), et permute \(q_+\leftrightarrow q_-\). Le couple est ainsi la lecture spectrale d'un même antécédent angulaire, et non deux paramètres ajustés séparément.

\begin{definition}[Réponse constitutive]
\begin{equation}
\mathcal V_{90,120}=(I-T^{90})(I-T^{120})^{-1}.
\label{eq:vacuum-operator}
\end{equation}
\end{definition}

Comme \(\|T\|<1\), le dénominateur est positif et inversible. Le calcul fonctionnel donne

\begin{equation}
\mathcal V_{90,120}=r_+P_+ +r_-P_-,
\qquad
r_\pm=\frac{1-q_\pm^{90}}{1-q_\pm^{120}}.
\label{eq:vacuum-eigenvalues}
\end{equation}

L'affirmation d'inversibilité mérite d'être explicitée. Pour chaque feuillet,
\(0<q_\pm^{120}<1\) ; par conséquent
\[
I-T^{120}=(1-q_+^{120})P_++(1-q_-^{120})P_-
\]
est strictement positif et
\[
(I-T^{120})^{-1}
=\frac{P_+}{1-q_+^{120}}+
 \frac{P_-}{1-q_-^{120}}.
\]
La réponse \(\mathcal V_{90,120}\) ne s'obtient pas en résolvant quatre
équations scalaires séparées : elle est le résultat du calcul fonctionnel d'une
seule contraction positive. Cette observation empêche que perméabilité,
permittivité, impédance et propagation perdent l'antécédent qu'elles partagent.

\begin{theorem}[Détermination spectrale conjointe des quatre relations constitutives]
Les quatre caractères constitutifs
\begin{equation}
\widehat\varepsilon=r_+^2,\qquad
\widehat\mu=r_-^2,\qquad
\widehat Z=\frac{r_-}{r_+},\qquad
\widehat c=\frac1{r_+r_-}
\label{eq:vacuum-four}
\end{equation}
sont des fonctions spectrales de \(\mathcal V_{90,120}\). Ils satisfont
\begin{equation}
\widehat\mu\widehat\varepsilon=\widehat c^{-2},
\qquad
\frac{\widehat\mu}{\widehat\varepsilon}=\widehat Z^2.
\label{eq:vacuum-relations}
\end{equation}
\end{theorem}

\begin{proof}
Les identités s'obtiennent en substituant \cref{eq:vacuum-four}. Aucune carte dimensionnelle ni aucune valeur extérieure n'intervient.
\end{proof}

\begin{corollary}[Reconstruction des deux feuillets]
Les quatre caractères ne contiennent pas quatre degrés de liberté. N'importe lequel
des couples \((\widehat\mu,\widehat\varepsilon)\),
\((\widehat Z,\widehat c)\) ou \((\mathcal E,\mathcal B)\) reconstruit les
deux valeurs propres positives :
\begin{align}
r_-&=\sqrt{\widehat\mu}
=\sqrt{\widehat Z/\widehat c},
&r_+&=\sqrt{\widehat\varepsilon}
=\frac1{\sqrt{\widehat Z\widehat c}},
\label{eq:vacuum-reconstruct-r}\\
\log r_-&=\frac{\mathcal E+\mathcal B}{2},
&\log r_+&=\frac{\mathcal E-\mathcal B}{2}.
\label{eq:vacuum-reconstruct-log}
\end{align}
En particulier, l'orientation se perd si l'on conserve uniquement le produit,
mais se retrouve en ajoutant le quotient.
\end{corollary}

\begin{proof}
La première ligne est la racine positive des identités
\cref{eq:vacuum-four,eq:vacuum-relations}. La seconde inverse le système
linéaire
\(\mathcal E=\log r_-+\log r_+\),
\(\mathcal B=\log r_--\log r_+\).
\end{proof}

Ce corollaire exprime une propriété centrale de HMT : produit et quotient sont
des fonctionnelles complémentaires d'un même état. Le produit détermine la
propagation commune, tandis que le quotient conserve l'orientation
relative des feuillets intérieur et extérieur.

Dans le certificat V2,

\begin{align*}
r_+&=0.9999996285482493631786952281\ldots,\\
r_-&=0.9997241371895183641315165853\ldots,\\
\widehat\varepsilon&=0.9999992570966367027604416155\ldots,\\
\widehat\mu&=0.9994483504793269350899815559\ldots,\\
\widehat Z&=0.9997245085389372154555543466\ldots,\\
\widehat c&=1.0002763104861575261063862689\ldots.
\end{align*}

Ces numéraux sont la reconnaissance terminale de \cref{eq:vacuum-operator} ; ils ne définissent pas ses exposants \(90,120\).

\section{Factorisation harmonique des secteurs \(90\) et \(120\)}

Le semi-groupe harmonique est

\[
\langle90,120\rangle=30\langle3,4\rangle.
\]

Si \(s=q^{30}\), alors

\begin{equation}
\frac{1-q^{90}}{1-q^{120}}
=\frac{1-s^3}{1-s^4}
=\frac{1+s+s^2}{1+s+s^2+s^3}.
\label{eq:three-four-completion}
\end{equation}

L'identité sépare également la réponse de son défaut. Si
\[
F_{3\to4}(s)=\frac{1+s+s^2}{1+s+s^2+s^3},
\qquad 0<s<1,
\]
alors
\begin{equation}
d_{3\to4}(s):=1-F_{3\to4}(s)
=\frac{s^3}{1+s+s^2+s^3}>0.
\label{eq:vacuum-defect34}
\end{equation}
Le terme supplémentaire est conservé exactement comme défaut
positif. En outre,
\begin{equation}
F_{3\to4}'(s)
=-\frac{s^2(3+2s+s^2)}{(1+s+s^2+s^3)^2}<0.
\label{eq:vacuum-monotonic34}
\end{equation}
Par conséquent, l'ordre des canaux angulaires détermine univoquement
l'ordre des réponses constitutives. L'attribution de
\(\widehat\mu\) et \(\widehat\varepsilon\) aux feuillets respectifs s'obtient
analytiquement, sans recourir à leurs développements décimaux.

Chaque feuillet du vide constitue la complétion exacte de trois à quatre
termes d'un même mode élémentaire \(30\). En conséquence, perméabilité
et permittivité s'interprètent comme deux fonctionnelles quadratiques d'une même
complétion évaluée sur des feuillets conjugués, et non comme des paramètres
constitutifs indépendants.

Comme la fonction \((1-q^{90})/(1-q^{120})\) décroît dans \(0<q<1\) et \(q_+<q_-\), on obtient

\[
0<r_-<r_+<1,\qquad
0<\widehat\mu<\widehat\varepsilon<1,\qquad
0<\widehat Z<1<\widehat c.
\]

Les limites de la fonctionnelle déterminent sa géométrie. Lorsque \(q\to0^+\),
\(F(q)\to1\) : les quatre lectures s'approchent de l'unité et le défaut de
mémoire disparaît. Lorsque \(q\to1^-\),
\[
F(q)\longrightarrow\frac{90}{120}=\frac34.
\]
La réponse prend donc ses valeurs dans l'intervalle positif \((3/4,1)\). La
borne \(3/4\) n'est pas une constante ajustée : elle est le quotient des deux
longueurs de canal et la réalisation continue de la complétion
trois--à--quatre.

\section{Involution des feuillets et conservation de l'orientation}

L'échange de feuillets \(C^*\mapsto-C^*\) agit par

\begin{equation}
\widehat\mu\longleftrightarrow\widehat\varepsilon,\qquad
\widehat Z\longmapsto\widehat Z^{-1},
\qquad
\widehat c\longmapsto\widehat c.
\label{eq:vacuum-sheet}
\end{equation}

La propagation commune est paire ; l'impédance conserve l'orientation.
Cette distinction réapparaîtra dans Barbero et CKM : toutes les fonctionnelles
associées au changement de feuillet n'ont pas la même parité.

\section{Caractères logarithmiques et compatibilité traciale}

Définissons

\[
\mathcal E=\log(r_+r_-),\qquad
\mathcal B=\log(r_-/r_+).
\]

Avec la trace normalisée \(\tau=\Tr/2\),

\begin{equation}
2\tau(\log\mathcal V)=\mathcal E,
\qquad
-2\tau(R\log\mathcal V)=\mathcal B.
\label{eq:vacuum-traces}
\end{equation}

Le couple \((\mathcal E,\mathcal B)\) constitue un système de coordonnées linéaire du logarithme
opératoriel :
\begin{equation}
\log\mathcal V
=\frac{\mathcal E}{2}I-\frac{\mathcal B}{2}R.
\label{eq:vacuum-log-decomposition}
\end{equation}
La composante scalaire est invariante sous l'échange de feuillets ; la
composante orientée change de signe. Cette décomposition démontre
directement la parité de la propagation et le caractère orienté de
l'impédance.

L'amplification nonadique

\[
\mathcal V_m=\mathcal V\otimes I_{9^m}
\]

est compatible avec les espérances \(E_m=\mathrm{id}\otimes\tau_9\) :

\[
E_m(\mathcal V_{m+1})=\mathcal V_m.
\]

Par conséquent, \(\mathcal E\) et \(\mathcal B\) ne sont pas des particularités d'une matrice \(2\times2\) ; ils subsistent dans la tour UHF et dans sa clôture traciale \(II_1\).

Plus précisément, pour tout polynôme \(p\) puis, par continuité,
pour toute fonction continue sur le spectre,
\begin{equation}
E_m\bigl(p(\mathcal V_{m+1})\bigr)=p(\mathcal V_m).
\label{eq:vacuum-functional-refinement}
\end{equation}
La compatibilité ne préserve pas seulement deux nombres : elle préserve toute
l'algèbre fonctionnelle engendrée par la réponse. En particulier, elle conserve ses
projecteurs spectraux, ses puissances, son logarithme et les formes quadratiques
qui mesurent les défauts \cref{eq:vacuum-defect34}.

\section{Réalisation dimensionnelle des relations constitutives}

La mécanique dimensionnelle détermine

\begin{align}
Z_{\rm vac}&=\widehat Z\,u_Su_Q^{-2},\\
c_{\rm vac}&=\widehat c\,u_C,\\
\mu_{\rm vac}&=\widehat\mu\,u_Su_C^{-1}u_Q^{-2},\\
\varepsilon_{\rm vac}&=\widehat\varepsilon\,u_S^{-1}u_C^{-1}u_Q^2.
\label{eq:vacuum-sections}
\end{align}

Les identités constitutives se maintiennent dans les droites :

\begin{equation}
\mu_{\rm vac}\varepsilon_{\rm vac}=c_{\rm vac}^{-2},
\qquad
\mu_{\rm vac}/\varepsilon_{\rm vac}=Z_{\rm vac}^2.
\end{equation}

La carte SI est postérieure. Il n'est pas nécessaire de déclarer \(c\) primitif et de résoudre ensuite \(\mu\varepsilon=c^{-2}\) : les trois objets descendent conjointement de \(\mathcal V\) et des droites déclarées.

La séparation des niveaux s'exprime par le diagramme
\[
\begin{array}{c}
(A,C^*)\longrightarrow T\longrightarrow\mathcal V
\longrightarrow
(\widehat\mu,\widehat\varepsilon,\widehat Z,\widehat c)
\\[2mm]
\Big\downarrow\text{ cartes de droite}\hspace{36mm}
\Big\downarrow\text{ évaluation finale}
\\[2mm]
(u_S,u_Q,u_C)\longrightarrow
(\mu_{\rm vac},\varepsilon_{\rm vac},Z_{\rm vac},c_{\rm vac})
\longrightarrow\text{ coordonnées SI}.
\end{array}
\]
La première ligne est adimensionnelle et opératorielle ; la deuxième détermine les types de
grandeur ; la troisième sélectionne une carte métrologique. Modifier la carte finale
ne modifie ni les valeurs propres ni les identités constitutives. En ce
sens précis, l'unification mathématique--physique consiste à distinguer
l'invariant généalogique de son système de coordonnées, et non à supprimer la
structure dimensionnelle.

\begin{proposition}[Minimalité spectrale]
Soit \(W\) un opérateur positif à deux feuillets de spectre simple
\(\{r_+,r_-\}\). Si quatre caractères positifs satisfont
\(\mu\varepsilon=c^{-2}\) et \(\mu/\varepsilon=Z^2\), si la propagation
commune est déterminée par \(c^{-1}=\det W=r_-r_+\), et si son orientation
fixe \(Z=r_-/r_+\), alors nécessairement
\[
(\mu,\varepsilon,Z,c)
=\left(r_-^2,r_+^2,\frac{r_-}{r_+},\frac1{r_-r_+}\right).
\]
Par conséquent, \cref{eq:vacuum-four} n'est pas une paramétrisation parmi d'autres :
c'est la réalisation positive minimale compatible avec le produit, le quotient et
l'orientation.
\end{proposition}

\begin{proof}
De \(Z=r_-/r_+\) et \(c^{-1}=\det W=r_-r_+\), on obtient de manière unique les
racines positives. Les deux identités constitutives imposent ensuite
\(\mu=r_-^2\) et \(\varepsilon=r_+^2\).
\end{proof}

